% ID: SEL-D0084-OB-E01
% ESTADO: cotejada. Segunda pasada visual y borrador compilado revisados.
% Fuente DOC: Selectividad/sel_2012_matematicas_aplicadas/6 MATEMÁT. APLIC. CC. SS/EXAMEN 6-ANDALUCÍA.doc
% SHA-256 original: 61f8d5d809881c49e9e9212281f7028516338e026524e55d572ab691ec5f416d
% Página convertida: 2; opción B; ejercicio 1.

\textbf{(2.5 puntos)} En una carpintería se construyen dos tipos de estanterías: grandes y pequeñas, y se tienen para ello $60\,m^2$ de tableros de madera. Las grandes necesitan $4\,m^2$ de tablero y las pequeñas $3\,m^2$. El carpintero debe hacer como mínimo 3 estanterías grandes, y el número de pequeñas que haga debe ser, al menos, el doble del número de las grandes. Si la ganancia por cada estantería grande es de 60 euros y por cada una de las pequeñas es de 40 euros, \textquestiondown{}cuántas debe fabricar de cada tipo para obtener el máximo beneficio?
